\subsection{Version 4: two editions from shared mathematical sources}
The complete v3 PDF is preserved as
\artifact{paper/versions/kourovka-experiment-2026-09-19-v3.pdf}.
The following literal source excerpts identify the version and dependency
changes. Appendix~\ref{app:two-editions} describes the two presentations.

\paragraph{Version identity}

\noindent\textbf{Version 3.}\begin{quote}\small
Version 3 --- 19 September 2026
\end{quote}

\noindent\textbf{Version 4.}\begin{quote}\small
Version 4 --- 19 September 2026
\end{quote}

Both entry points now use the shared version identity; the v3 PDF is preserved.

\paragraph{A shared mathematical dependency}

\noindent\textbf{Version 3.}\begin{quote}\small
The cyclic-vector argument of Section~\ref{cand:10.35}
\end{quote}

\noindent\textbf{Version 4.}\begin{quote}\small
The cyclic-vector argument of Lemma~\ref{lem:matrix-centralizer}
\end{quote}

The general centralizer fact has its own shared lemma, so omitting the withdrawn application removes no premise of 20.90.

\paragraph{Attribution in both editions}

\noindent\textbf{Version 3.}\begin{quote}\small
Version 2 attribution.
\end{quote}

\noindent\textbf{Version 4.}\begin{quote}\small
Related work.
\end{quote}

The existing attribution and its mathematical content remain shared; the heading no longer requires the reader to follow the revision history.

The original question and author blocks are additions, drawn from the
shared metadata. Three representative arguments were extracted into
separate input files; their mathematical text is preserved. The same
partial and prior-work arguments appear in both editions, while the
mathematical edition omits the experimental narrative, review correspondence,
bounded unsuccessful searches and withdrawn 10.35 application. The current
index is generated separately from the frozen deadline ledger. Preservation
checks and the source changes are recorded in
\artifact{paper/reviews/v4-2026-09-19/}.
